\subsection{Economic Incentives in AI Development}
\label{sec:8.3.1}
Financial incentive mostly determines what AI development produces. Google, Meta, and Amazon commit enormous resources to AI research aimed at operational efficiency, deeper user engagement, and market dominance, and that spending drives technical progress — successive leaps in natural-language processing capability among them \autocite{maslej2025index}. What that spending does not automatically produce is a system whose outputs serve the people using it rather than the shareholders funding it.

A recurring case is the profit-driven AI system that optimizes for engagement at the expense of the people engaged. Recommendation algorithms on social media platforms have spread misinformation and hardened political polarization, because a system tuned to maximize time-on-platform has no built-in reason to prefer accurate content over inflammatory content \autocite{germano2026ranking}.

The same economic logic that rewards engagement over accuracy also rewards scale over access: the firms with the deepest resources for AI research pull further ahead of smaller competitors, independent researchers, and everyone else who lacks comparable capital.

None of that is a fixed property of the technology. It is a property of who is paying, which is why the source of research funding is a policy instrument and not an accounting detail.
